\section{The tightening operator and what limits it}\label{sec:foundations}

This section fixes the operator studied in the rest of the paper, records
the order properties on which all later results rest, and identifies three
limits on further tightening. The order arguments are
elementary and not new. Monotone iterations of optimality-based domain
reduction, their limits, and their independence of the variable order are
studied by Caprara and
Locatelli~\cite{caprara2010-global-optimization-problems-and-domain}, and
greatest fixed points of feasibility-based tightening by Belotti et
al.~\cite{belotti2010-feasibility-based-bounds-tightening-via-fixed-points,belotti2012fbbt}. We state the facts in the form used later and
make each hypothesis explicit, because several later conclusions fail when
one of them is dropped.

\subsection{Problem, relaxation family, and order hypotheses}

Consider
\begin{equation}\label{eq:problem}
  f^*=\min\{f(x):x\in X\cap B_0\},
\end{equation}
where $B_0=[\ell^0,u^0]\subset\R^n$ is a bounded box and the closed set $X$
collects all other constraints, including integrality. We assume that
$f^*$ is attained and write $x^*$ for a global minimizer. A box is a set
$B=[\ell,u]=\{x:\ell\le x\le u\}$ with $\ell\le u$; degenerate boxes are
allowed, and the empty set is treated separately. We write
\[
  w(B)=\max_i(u_i-\ell_i),\qquad p(B)=(-\ell,u)\in\R^{2n}.
\]
For nonempty boxes, $B'\subseteq B$ holds exactly when $p(B')\le p(B)$
componentwise, so shrinking a box decreases its endpoint vector. When
tightening is applied at a node of a branch-and-bound tree, $B_0$ is the
node box and~\eqref{eq:problem} is the node problem.

A \emph{relaxation family} assigns to every box $B$ in a collection
$\mathfrak B$ a relaxation with projected objective
$\phi_B:B\to\R\cup\{+\infty\}$. The collection $\mathfrak B$ contains all
subboxes of $B_0$; Section~\ref{sec:local-rates} also uses boxes around a
minimizer that may extend beyond $B_0$, and then $\mathfrak B$ consists of
all subboxes of a larger domain box (for McCormick relaxations of
polynomials, all boxes in $\R^n$). A lifted relaxation has a lifted set
$R(B)\subseteq\R^n\times\R^m$ and objective $v$, and
\begin{equation}\label{eq:projected}
  \phi_B(x)=\inf\{v(x,y):(x,y)\in R(B)\},\qquad\inf\emptyset=+\infty.
\end{equation}
Points that violate the relaxed constraints therefore have projected value
$+\infty$. The two order hypotheses are
\begin{align}
  \phi_B(x)&\le f(x)&&(x\in X\cap B),\label{eq:validity}\\
  \phi_{B'}(x)&\ge\phi_B(x)&&(B'\subseteq B,\ x\in B').
  \label{eq:monotonicity}
\end{align}
We call~\eqref{eq:validity} validity and~\eqref{eq:monotonicity}
monotonicity of the family. In addition, we assume throughout that each
$\phi_B$ is lower semicontinuous, so that its sublevel sets in $B$ are
compact. When $\phi_B$ is given by~\eqref{eq:projected}, we assume that
$R(B)$ is compact and $v$ is continuous. Then the infimum
in~\eqref{eq:projected} is attained whenever it is finite, $\phi_B$ is
lower semicontinuous, and every projected sublevel set is the projection
of the corresponding lifted sublevel set. Without attainment, an
implementation that optimizes over the lifted cutoff set can compute a
smaller set than the projected one, and the two must not be identified.

Monotonicity is a hypothesis on the \emph{whole} relaxation construction:
every row, auxiliary bound, retained cut, and the objective. It does not
follow from the name of the relaxation. A sufficient condition is
\emph{lifted monotonicity}: $B'\subseteq B$ implies $R(B')\subseteq R(B)$,
with the same objective $v$. Lifted monotonicity is strictly stronger than
\eqref{eq:monotonicity}; it is what allows the same lifted points to be
reused on a larger box, and Section~\ref{sec:certificates} relies on it for
that purpose. The following constructions are monotone.
\begin{enumerate}[label=(\alph*)]
\item Termwise relaxations of quadratic objectives that keep convex
  separable terms exact and replace each bilinear term $x_ix_j$, $i\ne j$,
  by its McCormick envelope. The four McCormick inequalities describe the
  convex hull of the graph of $x_ix_j$ over the current rectangle. The graph
  over a subrectangle is a subset of the graph over the rectangle, so the
  convex hulls are nested.
\item Lifted linear relaxations with fixed linear rows and McCormick
  inequalities for lifted variables $y_k$ that stand for products $x_ix_j$,
  including squares $i=j$, for which
  the four inequalities reduce to the two endpoint tangents and the secant.
  Section~\ref{sec:certificates} uses this family as its exact reference
  family and proves its lifted monotonicity.
\item Composite McCormick relaxations of factorable functions in which
  every intermediate relaxation is clipped to its interval bounds; see
  Section~\ref{sec:local-rates-factorable}.
\end{enumerate}
Monotonicity can fail even when the construction is determined by the box
alone. Relax $x^2$ on $[\ell,u]$ by the tangents $y\ge2px-p^2$ at the five
equally spaced points $p=\ell+k(u-\ell)/4$, $k=0,\dots,4$, and by the secant
$y\le(\ell+u)x-\ell u$, with objective $v=y$; these are the square rows of the
numerical policy of Section~\ref{sec:algorithms-policy} without its retained
points. On the box the tangents lie below $x^2$ and hence below the secant,
so $\phi_{[\ell,u]}(x)=\max_k(2p_kx-p_k^2)$. The five points for $[-1,1]$
include $0$, whereas those for $[-1/2,1]$ are $-1/2,-1/8,1/4,5/8,1$. Hence
$\phi_{[-1,1]}(0)=0>-1/64=\phi_{[-1/2,1]}(0)$,
contrary to~\eqref{eq:monotonicity}. The operator~\eqref{eq:operator}
defined below is affected as well: at $U=1/4$ it maps $[-1,1]$ to
$[-1/2,1/2]$ but the smaller box $[-1/2,1]$ to $[-1/2,41/80]$, because the
point $(x,y)=(41/80,1/4)$ satisfies every row on $[-1/2,1]$ but violates the
tangent at $1/2$ of $[-1,1]$. Tangents at the two endpoints alone, as
in~(b), do not have this defect. Constructions that depend on more than the
box need not be monotone either. Cutting planes separated at earlier
relaxation solutions and kept in a pool depend on the search history, and
rows dropped to save time, or a relaxation that is solved only
approximately, can produce a larger feasible set on a smaller box. For such
constructions the statements below require a separate argument.

The relaxation bound on a box is
\[
  L(B)=\inf_{x\in B}\phi_B(x).
\]
Monotonicity gives $L(B')\ge L(B)$ whenever $B'\subseteq B$: for $x\in B'$
one has $\phi_{B'}(x)\ge\phi_B(x)\ge L(B)$.

\subsection{The operator and update schemes}

For a cutoff $U$, which in practice is the value of an incumbent, define
\begin{equation}\label{eq:operator}
  K_U(B)=\{x\in B:\phi_B(x)\le U\},\qquad
  T_U(B)=\hull K_U(B),
\end{equation}
where $\hull S$ is the smallest closed box containing $S$, and
$\hull\emptyset=\emptyset$. We also set $K_U(\emptyset)=T_U(\emptyset)=\emptyset$,
so that repeated rounds remain defined after an empty cutoff set; the
inclusions below then hold with $\emptyset$ as the smallest element. If
$K_U(B)\ne\emptyset$, then
$T_U(B)=[\ell',u']$ with
\[
  \ell'_i=\min\{x_i:x\in K_U(B)\},\qquad
  u'_i=\max\{x_i:x\in K_U(B)\},\qquad i=1,\dots,n.
\]
These $2n$ support problems are convex programs when $\phi_B$ is convex,
and linear programs for lifted polyhedral relaxations. The cutoff slack of
$U$ is $\epsilon=U-f^*$. It is nonnegative when $U$ is the value of a point
of $X\cap B_0$; at a node, a global incumbent outside the node box can have
a value below the node optimum, and then the node can be pruned. Results
that need a nonempty limit or a minimizer in every iterate assume
$U\ge f^*$ explicitly.

Solvers organize the support problems in different ways, and the
difference matters for exact rate statements.
\begin{itemize}
\item A \emph{Jacobi round} solves all $2n$ support problems on the one
  frozen relaxation built on $B$ and returns $T_U(B)$.
\item A \emph{directional update} replaces one endpoint of the current box
  by the corresponding support over the cutoff set of the relaxation built
  on the current box. If that set is empty, the update returns $\emptyset$,
  and every later update of $\emptyset$ returns $\emptyset$.
\item A \emph{sequential round} applies $2n$ directional updates, one per
  signed coordinate direction, rebuilding the relaxation after each update.
  We write $S_U(B)$ for its result; it depends on the order of directions.
\item A \emph{selective round} applies directional updates to a subset of
  the directions.
\end{itemize}

\begin{lemma}[Containment and order]\label{lem:order}
For all boxes $B',B$ and cutoffs $U'\le U$:
\begin{enumerate}[label=(\alph*)]
\item $\{x\in X\cap B:f(x)\le U\}\subseteq K_U(B)\subseteq T_U(B)\subseteq B$;
\item if $B'\subseteq B$, then $K_U(B')\subseteq K_U(B)$ and
  $T_U(B')\subseteq T_U(B)$;
\item $T_{U'}(B)\subseteq T_U(B)$.
\end{enumerate}
\end{lemma}

\begin{proof}
(a) For a feasible $x\in B$ with $f(x)\le U$, validity gives
$\phi_B(x)\le f(x)\le U$. The remaining inclusions hold because
$K_U(B)\subseteq B$ and $B$ is a closed box. (b) For $x\in K_U(B')$,
monotonicity gives $\phi_B(x)\le\phi_{B'}(x)\le U$, so $x\in K_U(B)$; the
box hull preserves inclusion. (c) $K_{U'}(B)\subseteq K_U(B)$.
\end{proof}

\begin{lemma}[Update schedules]\label{lem:sequential}
Let $C_0=B$ and let $C_1,C_2,\dots$ be the boxes produced by a sequence of
directional updates at cutoff $U$.
\begin{enumerate}[label=(\alph*)]
\item $T_U^m(B)\subseteq C_m$ for every $m$.
\item If every one of the $2n$ directions is updated at least once among
  the updates $j+1,\dots,m$, then $C_m\subseteq T_U(C_j)$. In particular,
  every sequential round satisfies $T_U^{2n}(B)\subseteq S_U(B)\subseteq T_U(B)$.
\item If every direction is updated infinitely often, then
  $\bigcap_mC_m=\bigcap_kT_U^k(B)$.
\end{enumerate}
\end{lemma}

\begin{proof}
(a) A directional update of a box $C$ changes one endpoint to the
corresponding endpoint of $T_U(C)$ and keeps the others, so its result
contains $T_U(C)$. Induction with Lemma~\ref{lem:order}(b) gives
$C_{m}\supseteq T_U(C_{m-1})\supseteq T_U(T_U^{m-1}(B))=T_U^m(B)$.
(b) Each update after step $j$ optimizes over a cutoff set contained in
$K_U(C_j)$, by Lemma~\ref{lem:order}(b). An endpoint updated in this way
is therefore at least as tight as the corresponding endpoint of
$T_U(C_j)$, and later updates never move it outward. Every endpoint is
updated at least once, which proves $C_m\subseteq T_U(C_j)$; part (a) with
$m=2n$ gives the other inclusion for a sequential round. (c) Partition the
schedule into consecutive finite blocks, each containing every direction
at least once, and let $t_j$ be the end of block $j$, with $t_0=0$. By (b)
and monotonicity of $T_U$, $C_{t_j}\subseteq T_U(C_{t_{j-1}})\subseteq
T_U^j(B)$, and by (a), $T_U^{t_j}(B)\subseteq C_{t_j}$. Each block is
nonempty, so $t_j\ge j$. The sequences $(T_U^{t_j}(B))_j$ and $(T_U^j(B))_j$
are therefore subsequences of the nested sequence $(T_U^k(B))_k$ with
indices tending to infinity, and both have the intersection
$\bigcap_kT_U^k(B)$. The two inclusions give the same intersection for
$(C_{t_j})_j$, and this equals $\bigcap_mC_m$ because the boxes $C_m$ are
nested.
\end{proof}

Part (c) is the order independence of the limit studied by Caprara and
Locatelli~\cite{caprara2010-global-optimization-problems-and-domain},
stated here for every fair schedule of directional updates and for Jacobi
rounds; it uses only monotonicity. The lemma also determines which bounds transfer between
update schemes. Every upper bound on the box produced by a Jacobi round
also bounds the box produced by a complete sequential round, so upper
bounds on widths and rates transfer. A selective round that leaves some
directions unprocessed need not be contained in $T_U(B)$ and inherits no
upper bound, but after $m$ updates it still contains $T_U^m(B)$, so lower
enclosures transfer with the number of updates. Lower bounds on the box
after a complete round do not transfer from Jacobi to sequential rounds,
because a sequential round can be strictly faster. For example, let $n=1$,
$f(x)=x^2$, $U=0$, and $\phi_{[\ell,u]}(x)=x^2-(u-\ell)^2/16$, a valid,
convex, and monotone family. A Jacobi round maps $[-h,h]$ to $[-h/2,h/2]$.
A sequential round that updates the lower and then the upper endpoint maps
$[-a,b]$ to $[-a',b']$ with $a'=(a+b)/4$ and $b'=(a'+b)/4=(a+5b)/16$,
provided $b\le3a$ and $a\le11b$; the ratio $b/a$ stays in $[1/2,1]$ along
the iteration from $a=b$, so these conditions hold throughout. The
sequential widths therefore decrease like $\lambda_s^k$ with
$\lambda_s=(9+\sqrt{17})/32\approx0.410$, the larger eigenvalue of
$\left[\begin{smallmatrix}1/4&1/4\\1/16&5/16\end{smallmatrix}\right]$,
whereas Jacobi widths decrease exactly like $2^{-k}$.

\subsection{Iterates, fixed boxes, and the limit}

Let $B_{k+1}=T_U(B_k)$, starting from $B_0$, and let
$B_\infty=\bigcap_kB_k$. By Lemma~\ref{lem:order} the iterates are nested.
If $U\ge f^*$, every iterate contains every global minimizer, so $B_\infty$
is a nonempty compact box. The following lemma describes the part of the
box that no valid tightening at cutoff $U$ can remove.

\begin{lemma}[Hull of the original sublevel set]\label{lem:sublevel-hull}
Let $H_U=\hull\{x\in X\cap B_0:f(x)\le U\}$. Every iterate contains $H_U$,
and $T_U(H_U)=H_U$ if $H_U\ne\emptyset$. In particular, if $x^*$ and
$x^{**}$ are two global minimizers, then
$H=\hull\{x^*,x^{**}\}\subseteq B_\infty$ for every $U\ge f^*$ and
\[
  L(B_\infty)\le L(H).
\]
\end{lemma}

\begin{proof}
By Lemma~\ref{lem:order}(a) and induction, every iterate contains each
feasible point with objective at most $U$, hence their hull, because the
iterates are closed boxes. The same points lie in $K_U(H_U)$, so
$H_U\subseteq\hull K_U(H_U)=T_U(H_U)\subseteq H_U$. Finally
$H\subseteq H_{f^*}\subseteq H_U\subseteq B_\infty$, and monotonicity of $L$
gives $L(B_\infty)\le L(H)$.
\end{proof}

Caprara, Locatelli, and Monaci call $H_U$ the lower limit of
optimality-based domain
reduction~\cite{caprara2016-theoretical-and-computational-results-about}.
They identify a two-variable class for which iterated reduction reaches it
and give examples in which reduction does not move any bound although
$H_U$ is a single point. Lemma~\ref{lem:sublevel-hull} also limits iterated
tightening at the root of a symmetric model: if the problem has several
global minimizers related by a symmetry, the limit box contains their hull
for every $U\ge f^*$, so the iterates cannot converge to a single point, and
the relaxation bound after tightening is at most the relaxation bound on
that hull. Tightening can still remove parts of the box outside the hull.

The box $H_U$ depends only on the original problem and the cutoff. The
relaxation can keep more than $H_U$.

\begin{definition}[Fixed box]
A nonempty box $P$ is a \emph{fixed box} of $T_U$ if $T_U(P)=P$.
\end{definition}

Since $T_U(P)\subseteq P$ always, $P$ is fixed exactly when
$P\subseteq T_U(P)$, that is, when every face of $P$ is attained by a point
of $K_U(P)$.

\begin{lemma}[Persistence of fixed boxes]\label{lem:fixed-persist}
Let $P$ be a fixed box of $T_U$ and $P\subseteq B$. Then $P\subseteq T_{U'}(B)$
for every $U'\ge U$, and $P$ is contained in the result of every
directional update of $B$ at a cutoff $U'\ge U$. Consequently every Jacobi,
sequential, or selective iteration started at a box containing $P$, with
cutoffs at least $U$, keeps $P$ in all its iterates.
\end{lemma}

\begin{proof}
By Lemma~\ref{lem:order}, $K_U(P)\subseteq K_U(B)\subseteq K_{U'}(B)$. The
point of $K_U(P)$ that attains the lower face $a_i$ of $P$ in coordinate $i$
lies in $K_{U'}(B)$, so the minimum of $x_i$ over $K_{U'}(B)$ is at most
$a_i$; the upper faces are analogous. Hence no directional update, and no
Jacobi round, moves an endpoint of $B$ past the corresponding face of $P$.
Induction over the updates proves the last claim.
\end{proof}

Lemma~\ref{lem:fixed-persist} is the order-theoretic core of the
certificates in Section~\ref{sec:certificates}: finitely many points of
$K_U(P)$ that attain all faces of $P$ prove that $P$ is fixed, and hence
bound the total movement of every endpoint under all later rounds of the
same relaxation family at cutoffs at least $U$.
Section~\ref{sec:certificates} lowers this cutoff to the largest relaxation
value of the points. Fixed boxes can also be combined.

\begin{lemma}[Greatest fixed box]\label{lem:greatest-fixed}
Let $\mathcal P$ be a nonempty family of fixed boxes of $T_U$ contained in
$B_0$. Then $\hull\bigcup_{P\in\mathcal P}P$ is a fixed box. Consequently,
the hull $P_U$ of all fixed boxes in $B_0$, with $P_U=\emptyset$ if there is
none, contains every fixed box in $B_0$ and, when nonempty, is itself the
greatest fixed box in $B_0$. Every Jacobi iterate contains $P_U$, and
$H_U\subseteq P_U$.
\end{lemma}

\begin{proof}
The hull $Q=\hull\bigcup_{P\in\mathcal P}P$ is a nonempty box in $B_0$. Each
$P\in\mathcal P$ lies in $Q$, so Lemma~\ref{lem:order}(b) gives
$P=T_U(P)\subseteq T_U(Q)$. The closed box $T_U(Q)$ therefore contains $Q$,
and Lemma~\ref{lem:order}(a) gives $T_U(Q)=Q$. Applied to the family of all
fixed boxes in $B_0$, when it is nonempty, this shows that $P_U$ is a fixed
box; by construction it contains every fixed box in $B_0$.
Lemma~\ref{lem:fixed-persist} places $P_U$ in every Jacobi iterate; for
$P_U=\emptyset$ there is nothing to prove. If $H_U\ne\emptyset$, it is a
fixed box in $B_0$ by Lemma~\ref{lem:sublevel-hull}, so $H_U\subseteq P_U$.
\end{proof}

Proposition~\ref{prop:join} in Section~\ref{sec:certificates} is the
finite-certificate form of the first statement. In lattice terms, the
subboxes of $B_0$ together with $\emptyset$ form a complete lattice under
inclusion, $T_U$ is monotone on it, and $P_U$ is its greatest fixed point;
the lemma is the elementary case of Tarski's
fixed-point theorem~\cite{tarski1955fixpoint} that we need. It uses
monotonicity alone. Identifying $P_U$ with the limit of the iteration needs
a closedness condition.

\begin{proposition}[The limit and the greatest fixed box]\label{prop:fixed-limit}
Suppose every Jacobi iterate is nonempty, for example because $U\ge f^*$.
The limit $B_\infty$ of the Jacobi iterates is a nonempty compact box. It
contains $P_U$, and hence $H_U$ and every fixed box in $B_0$, and it is
also the limit of every sequence of directional updates that updates each
direction infinitely often. Suppose
in addition that the family is closed along decreasing sequences in the
following sense:
\begin{equation}\label{eq:closed-family}
  C_k\downarrow C,\quad x_k\in K_U(C_k),\quad x_k\to x
  \quad\Longrightarrow\quad x\in K_U(C),
\end{equation}
where $C_k\downarrow C$ means that the boxes are nested with intersection
$C$. Then $B_\infty$ is a fixed box, and hence $B_\infty=P_U$.
\end{proposition}

\begin{proof}
The iterates are nested nonempty compact boxes, so their intersection is a
nonempty compact box. The containments are
Lemma~\ref{lem:greatest-fixed}, and the statement about schedules is
Lemma~\ref{lem:sequential}(c).

Assume~\eqref{eq:closed-family}. For each $k$ and each direction, say the
upper endpoint of coordinate $i$, the endpoint of $B_{k+1}$ is
$\max\{x_i:x\in K_U(B_k)\}$, attained at some $y^k\in K_U(B_k)$ by
compactness. A subsequence of $(y^k)$ converges to some $\bar y$, and
\eqref{eq:closed-family} gives $\bar y\in K_U(B_\infty)$. Its coordinate
$\bar y_i$ is the limit of the upper endpoints of $B_{k+1}$, which is the
upper endpoint of $B_\infty$ in coordinate $i$. Repeating the argument for
every direction shows that every face of $B_\infty$ is attained in
$K_U(B_\infty)$, so $B_\infty$ is a fixed box in $B_0$. Hence
$B_\infty\subseteq P_U$, and the first part gives equality.
\end{proof}

The same argument shows that~\eqref{eq:closed-family} makes $T_U$
continuous along decreasing sequences of boxes,
$T_U(\bigcap_kC_k)=\bigcap_kT_U(C_k)$, which is the property that
identifies the limit of the iterates with the greatest fixed box.
Condition~\eqref{eq:closed-family} holds whenever the projected objective,
as a function $(x,\ell,u)\mapsto\phi_{[\ell,u]}(x)$ on
$\{\ell\le x\le u\}$, is lower semicontinuous: if $C_k\downarrow C$,
$x_k\in K_U(C_k)$, and $x_k\to x$, then $x\in C$, the endpoints of $C_k$
converge to those of $C$, and
$\phi_C(x)\le\liminf_k\phi_{C_k}(x_k)\le U$. The termwise relaxations
of~(a) are jointly continuous in this sense. For lifted families, suppose
$R(B)=\{z\in Z:\ g_j(z;\ell,u)\le0,\ j=1,\dots,m\}$
with the box rows included, a compact set $Z$, functions $g_j$ continuous
in $(z,\ell,u)$, and a continuous objective $v$. If $x_k\in K_U(C_k)$,
choose lifts $z_k\in R(C_k)$ with $v(z_k)\le U$. A subsequence converges
to some $z$, the endpoints of $C_k$ converge monotonically to those of $C$,
and continuity gives $z\in R(C)$ and $v(z)\le U$; its projection is the
limit $x$. McCormick rows have coefficients that are polynomial in the box
endpoints, so they satisfy this condition, as do tangent rows at points
that depend continuously on the endpoints. Closedness and monotonicity are
independent hypotheses: the five-point tangent family above satisfies
\eqref{eq:closed-family} but is not monotone, and the next example is
monotone but violates~\eqref{eq:closed-family}.

Without~\eqref{eq:closed-family}, the limit of iterated tightening need not
be a fixed box, and can be strictly larger than the greatest fixed box.

\begin{example}[A limit that is not a fixed box]\label{ex:not-fixed}
Let $n=1$, $X=\{0\}$, $f\equiv0$, $B_0=[0,1]$, and $U=0$. For a box
$B=[r,s]\subseteq B_0$ let $\phi_B(x)=0$ if $x\le g(s)$ and $\phi_B(x)=+\infty$
otherwise, where $g(s)=s/4$ for $s\le1/2$ and $g(s)=(2s+1)/4$ for $s>1/2$.
The function $g$ is nondecreasing with $g(s)\le s$ on $[0,1]$. Hence the
family is valid ($\phi_B(0)=0$ whenever $0\in B$), monotone, and has closed
sublevel sets, and $T_0([0,s])=[0,g(s)]$. From $s_0=1$ the iterates satisfy
$s_k-1/2=2^{-k-1}$, so $B_\infty=[0,1/2]$, whereas
$T_0(B_\infty)=[0,1/8]$. A box $[r,s]\subseteq B_0$ is fixed only if
$g(s)\ge s$, that is, $s=0$. Hence $P_U=H_U=\{0\}$. Continuing from
$B_\infty$ gives $[0,1/8],[0,1/32],\dots$, which reach $P_U$ only as the
limit of a second infinite sequence.
\end{example}

\subsection{Three limits on further tightening}\label{sec:three-limits}

Lemmas~\ref{lem:sublevel-hull}, \ref{lem:fixed-persist},
and~\ref{lem:greatest-fixed} and Proposition~\ref{prop:fixed-limit}
identify three things that limit how much
further rounds can remove from a current iterate $B_k$, for a fixed
relaxation family and cutoff:
\begin{enumerate}[label=(\roman*)]
\item the hull $H_U$ of the original sublevel set, which no valid tightening
  at cutoff $U$ removes and which only a better incumbent or branching
  shrinks;
\item the fixed boxes of $T_U$, which no number of further rounds of the
  same relaxation family removes, and which only a smaller cutoff, a
  stronger relaxation, or branching can shrink; their hull $P_U$, the
  greatest fixed box when one exists and empty otherwise, contains $H_U$
  and every fixed box in $B_0$, and equals the limit $B_\infty$
  under~\eqref{eq:closed-family} (if some iterate is empty, both are
  empty);
\item the rate at which the iterates approach $B_\infty$, which determines
  how many rounds are needed for a given part of the possible reduction.
\end{enumerate}
These objects satisfy $H_U\cup P\subseteq P_U\subseteq B_\infty\subseteq B_k$
for every fixed box $P$ in $B_0$, although $H_U$ and $P$ need not contain
each other. For example, with $f(x)=x^2$ on $[-1,1]$, $U=0$, and the valid
family $\phi_B\equiv0$, one has $H_U=\{0\}$ and the fixed box $\{1\}$, and
there $P_U=B_\infty=[-1,1]$. If $P=[a,b]$ is a fixed box inside
the current box $B_k=[\ell^k,u^k]$, the total future movement of each
endpoint satisfies
\[
  \ell_i^\infty-\ell_i^k\le a_i-\ell_i^k,\qquad
  u_i^k-u_i^\infty\le u_i^k-b_i.
\]
The first item is a property of the problem, the second of the relaxation
family, and the third of the iteration. Section~\ref{sec:local-rates}
gives sufficient conditions near a minimizer under which fixed boxes exist
at every sufficiently small scale, so that tightening stalls, and
sufficient conditions under which the iterates contract geometrically, with
a limit of width $O(\sqrt\epsilon)$; these conditions do not cover every
relaxation family.
Section~\ref{sec:certificates} turns the second item into finite
certificates, and Section~\ref{sec:residual} bounds the remaining movement
under a uniform sensitivity bound. Sections~\ref{sec:algorithms}
to~\ref{sec:experiments} examine how far this information supports
decisions about when to start, stop, and repeat tightening in a solver.
